\clearpage
\chapter*{2024 Septiembre}
\addcontentsline{toc}{chapter}{2024 Septiembre}

\begin{enunciado}
    \hypertarget{2024 S 1}{}
    \textbf{1. (2,5 puntos)} En \(\mathbb{R}\) se considera la topología \(T\) cuya base es
    \[
        B = \{[a,b) \mid a,b \in \mathbb{R}, a < b\}.
    \]
    Estudiar la continuidad de la aplicación
    \[
        f_k : (\mathbb{R}, T) \to (\mathbb{R}, T)
    \]
    definida como \(f_k(x) = kx\) según los distintos valores de \(k \in \mathbb{R}\).
\end{enunciado}

\begin{enunciado}
    \hypertarget{2024 S 2}{}
    \textbf{2. (2,5 puntos)} Dada \(f : (\mathbb{R}, T_u) \to \mathbb{R}\), definida por \(f(x) = |x|\), sea \(S\) la topología final de \(f\) en \(\mathbb{R}\).
    \begin{enumerate}[label=\alph*)]
        \item Estudiar si \((\mathbb{R}, S)\) es un espacio de Hausdorff.
        \item Estudiar si \((\mathbb{R}, S)\) verifica los axiomas I y II de numerabilidad.
    \end{enumerate}
\end{enunciado}

\begin{enunciado}
    \hypertarget{2024 S 3}{}
    \textbf{3. (2,5 puntos)} En el intervalo \(X = [-1, 1]\) de la recta real se considera la topología
    \[
        T = \{A \subset X \mid 0 \notin A \text{ o bien } (-1, 1) \subset A\}.
    \]
    Estudiar si el espacio topológico \((X, T)\) es compacto.
\end{enunciado}

\begin{enunciado}
    \hypertarget{2024 S 4}{}
    \textbf{4. (2,5 puntos)} En \(\mathbb{R}\) se considera la topología \(T\) engendrada por la familia de subconjuntos
    \[
        F = \{(a, \infty) \mid a \in \mathbb{R}\}.
    \]
    \begin{enumerate}[label=\alph*)]
        \item Estudiar si \((\mathbb{R}, T)\) es normal.
        \item Caracterizar las sucesiones convergentes en \((\mathbb{R}, T)\).
    \end{enumerate}
\end{enunciado}